\section*{ID6219: Physical Consequences (Short)}
\paragraph{Metadata.}
PiID: 3377024242916 | RTEID: RTE:P24149.7581:T2.0923 | UUID: 37964272-4609-5b19-90f6-6d24b06e4408

\paragraph{Canonical Statement.}
\# Physical consequences (short) - **Orbital phase** \textbackslash{}(e\textasciicircum{}\{i\textbackslash{}ell\textbackslash{}theta\}\textbackslash{}) means the two-particle relative state carries orbital angular momentum (vortex). Collapse localizes particle-2 onto that ring/phase pattern relative to the measured \textbackslash{}(x\_1\textbackslash{}). - **δ-coupling** acts at \textbackslash{}(r=0\textbackslash{}): if your ring breathes inward so amplitude reaches \textbackslash{}(r\textbackslash{}approx0\textbackslash{}) the coupling can strongly entangle or transfer energy (can produce bound or resonant states). - **Breathing** \textbackslash{}(\textbackslash{}sigma\_r(t)\textbackslash{}) gives an explicit time modulation (inward–outward frequency) and so injects sidebands in the relative-frequency spectrum — that’s your “interjected frequency spikes.” - **Measurement timing** and the phase of the breathing control whether the post-measurement state is trapped (resonant coherent state) or disperses.

\paragraph{Canonical Equation.}
\[
r=0
\]

\paragraph{Dictionary.}
Physical consequences (short) - Orbital phase \textbackslash{}(e\textasciicircum{}i \textbackslash{}) means the two-particle relative state carries orbital angular momentum (vortex) [ID6219]

\paragraph{Encyclopedia.}
Physical Consequences (Short) is an algebraic definition, written r=0. It is an algebraic definition expressing the quantity directly in terms of the variables on its right-hand side. Within the Trawin Zero Point registry it serves as one element of the framework's mathematical narrative.
